\section*{Chapter \ref{ch:iv:betting-and-market-making-games} --- Betting and Market-Making Games}

\begin{solution}{iq:iv:betting-and-market-making-games:1}
The payoff averages $2 \times (2 + 4 + 6)/6 = 4$, the price: fair. The variance is high (half the time you receive
nothing), which matters for sizing but not for the price.
\iqlookfor{the expectation over the six faces and a word on variance.}
\end{solution}

\begin{solution}{iq:iv:betting-and-market-making-games:2}
3 to 1 against implies $\tfrac14$. At 30\%: $0.3 \times 3 - 0.7 \times 1 = 0.2$ per unit staked, a 20\% edge; the
bookmaker's implied probability is below yours.
\iqlookfor{converting odds to probability and computing the edge.}
\end{solution}

\begin{solution}{iq:iv:betting-and-market-making-games:3}
The gap is 0 with probability $6/36$ and $k = 1, \dots, 5$ with probability $2(6-k)/36$, so its mean is $70/36 \approx
1.94$ and its variance $210/36 - (70/36)^2 \approx 2.05$, a standard deviation of about 1.43. A market of 1.5 at 2.5 is
centred and about a third of a standard deviation each side; the distribution is skewed (a gap of 0 or 1 four times in nine),
so a buyer at 2.5 needs a gap of 3 or more, which comes up a third of the time. Widen it if the interviewer can see the
dice.
\iqlookfor{the distribution written down, mean and spread computed, and a width tied to both.}
\end{solution}

\begin{solution}{iq:iv:betting-and-market-making-games:4}
Centre 5, standard deviation $\sqrt{10 \times 0.25} \approx 1.58$; quote 4.5 at 5.5, say, and say that the
distribution is symmetric, so there is no reason to skew.
\iqlookfor{the binomial mean and spread, and a symmetric quote.}
\end{solution}

\begin{solution}{iq:iv:betting-and-market-making-games:5}
$\P(\max \le k) = (k/6)^3$, so $\E[\max] = \sum_{k=1}^6 \P(\max \ge k) = \sum_k \big(1 - ((k-1)/6)^3\big) =
119/24 \approx 4.96$. A quick check: it must exceed the two-dice value $161/36 \approx 4.47$.
\iqlookfor{the tail-sum formula and a \omterm{def:iv:phone-and-technical-screens:sanity}{sanity check} against the two-dice case.}
\end{solution}

\begin{solution}{iq:iv:betting-and-market-making-games:6}
With net odds $b = 3$, $f^\ast = p - q/b = 0.35 - 0.65/3 = 2/15 \approx 0.133$: commit about 13\% of wealth each time. The
wager has a large edge ($0.35 \times 3 - 0.65 = 0.40$ per unit) but loses almost two times in three, which keeps the
fraction small.
\iqlookfor{the Kelly formula for net odds $b$ and a reading of why it is small.}
\end{solution}

\begin{solution}{iq:iv:betting-and-market-making-games:7}
Maximise $0.3\ln(1 + 2f) + 0.4\ln(1 - f)$ (the zero outcome contributes nothing): $0.6/(1 + 2f) = 0.4/(1 - f)$ gives
$f = 1/7 \approx 0.143$. The growth rate there is about 0.0137 per bet (1.4\%), lower than at slightly smaller or
larger fractions, as the chapter's code checks.
\iqlookfor{setting up the log-growth first-order condition with several outcomes.}
\end{solution}

\begin{solution}{iq:iv:betting-and-market-making-games:8}
By \cref{prop:iv:betting-and-market-making-games:halving}, $\tfrac12$ at full Kelly and $(\tfrac12)^3 = \tfrac18$ at half
Kelly; a simulation of geometric Brownian motion in the chapter's code agrees. Full Kelly maximises long-run growth
only if the edge is known exactly; it is not, and overestimating the edge by a factor of two turns full Kelly into
double Kelly, whose growth is zero. At half the Kelly fraction the growth rate is still three quarters of its maximum, and deep drawdowns become far
rarer.
\iqlookfor{the drawdown formula and the argument from uncertainty in the edge.}
\end{solution}

\begin{solution}{iq:iv:betting-and-market-making-games:9}
The count is hypergeometric: mean $10 \times 26/52 = 5$, variance $10 \times \tfrac12 \times \tfrac12 \times 42/51
= 35/17$, standard deviation about 1.43; quote 4.5 at 5.5. After three reds, the other seven come from 49 cards
with 23 red: fair value $3 + 7 \times 23/49 = 44/7 \approx 6.29$.
\iqlookfor{sampling without replacement, and conditioning by removing the seen cards.}
\end{solution}

\begin{solution}{iq:iv:betting-and-market-making-games:10}
Before any trade, 10.5. The buy says the seen die is at least 4, mean 5: fair value $5 + 7 = 12$. The second buy
carries no new information (same die, same rule), so fair value stays 12: move your market to about 11.5 at 12.5 once,
and do not keep walking it up.
\iqlookfor{updating on what the trade reveals, and recognising a repeated trade as no new information.}
\end{solution}

\begin{solution}{iq:iv:betting-and-market-making-games:11}
Maximising $(1 - \alpha)h(1 - h/20) - \alpha(50 - h)^2/100$ over $h$ with $\alpha = 0.15$ gives $h \approx 11.4$
and a profit of about 1.93 per arrival (with $\alpha = 0.05$, $h \approx 10.4$ and 3.96;
\cref{fig:iv:betting-and-market-making-games:width}). With $\alpha = 0.3$ every half-width below 50 loses money:
the informed flow costs more than the uninformed pays. Do not quote, or quote so wide that nobody trades, and find
out who the informed traders are before quoting again.
\iqlookfor{the trade-off between width and informed losses, and recognising when not to trade.}
\end{solution}

\begin{solution}{iq:iv:betting-and-market-making-games:12}
With both bidding their estimate minus $k$, the winner is the one with the larger noise and pays $V + \max(e_1, e_2) -
k$, so its average profit is $k - \E[\max(e_1, e_2)]$. For two uniform variables on $[-10, 10]$ the expected maximum
is $10/3$, so the winner breaks even at $k = 10/3 \approx 3.3$. Winning is information: it tells you your estimate
was the higher one, and so probably too high (the winner's curse). The calculation ignores the edges of $V$'s range.
\iqlookfor{conditioning on winning, and the expected maximum of the noise.}
\end{solution}

\begin{solution}{iq:iv:betting-and-market-making-games:13}
By \omterm{def:iv:brainteasers-and-logic:backward}{backward induction} on (reds left, blacks left): stop when the continuation value is not positive. For 10 and 10 the
value is $26\,995/16\,796 \approx 1.61$; for a full deck, about 2.62. It is positive because you can stop when ahead
and are never forced to take the negative tail: stopping at zero is always available, which makes the value the
expectation of a non-negative option.
\iqlookfor{a dynamic programme over the deck's state, and the option-value intuition.}
\end{solution}

\begin{solution}{iq:iv:betting-and-market-making-games:14}
In units of 100, start at 10, move $\pm1$ with $p = 0.55$ for 100 rounds, absorbed at 0. The exact distribution
gives a ruin probability of about 0.090 and an expected final bankroll of about 19.55 units, 1\,955. Betting a
fixed fraction of wealth instead can never go to zero in finitely many rounds and scales the bet down after losses,
which is why fractional sizing survives streaks that ruin fixed bets; its median growth is set by the Kelly
analysis.
\iqlookfor{a small dynamic programme for ruin with a horizon, and the contrast with proportional betting.}
\end{solution}
